%% ma: L12-B08-0076
%% lop: 12
%% chuong: 2
%% bai: 8
%% dang: 12.8.8
%% loai: TLN
%% muc: TH
%% dap_an: "1,25"
%% nguon: "LaTeX giáo viên gửi 21/10/2026 (M05, Vectơ và hệ trục tọa độ trong không gian, 107 câu), Phần III Câu 16"
%% kien_thuc: "Bài 6 (quy tắc ba điểm), Bài 8 (biểu thức tọa độ các phép toán vectơ)"
%% trang_thai: chinh-thuc
%% ghi_chu: "Trùng ý tưởng với dạng 12.8.8 (hệ số âm, ba điểm) nên nhập cùng dạng; bỏ tiêu đề '(Tâm tỉ cự)'"
\begin{ex}
Trong không gian $Oxyz$, cho ba điểm $A(1;0;0)$, $B(0;2;0)$, $C(0;0;3)$. Gọi $I(a;b;c)$ là điểm thỏa mãn $2\overrightarrow{IA}+3\overrightarrow{IB}-\overrightarrow{IC}=\overrightarrow{0}$. Tính $S=a+b+c$.
\shortans{1,25}
\loigiai{$4\overrightarrow{OI}=2\overrightarrow{OA}+3\overrightarrow{OB}-\overrightarrow{OC}=(2;6;-3)$ nên $I\left(\dfrac12;\dfrac32;-\dfrac34\right)$, $S=1{,}25$.}
\end{ex}
